\documentclass[10pt,a4paper]{amsart}
\usepackage[margin=19mm]{geometry}
\usepackage{amsmath,amssymb,mathtools}
\usepackage[scr=rsfs,cal=euler]{mathalfa}
\usepackage{xfrac}
\usepackage[unicode,hidelinks,bookmarks=false]{hyperref}
\newcommand{\ie}{i.e.\ }
\newcommand{\cf}{cf.\ }
\newcommand{\resp}{respectively\ }
\newcommand{\Subsection}{\subsection}
\providecommand{\nobreakdash}{\nobreak\hbox{-}}
\setlength{\emergencystretch}{2em}
\multlinegap0pt
\usepackage[shortlabels]{enumitem}
\usepackage{amssymb}
\usepackage{mathtools}
\newtheorem{lemma}{Lemma}
\newcommand{\cA}{\mathcal A}
\newcommand{\cB}{\mathcal B}
\title{A new bound for the Kalton--Roberts constant}
\author{Boon Suan Ho, Tomasz Kania}
\date{Source: arXiv:2606.06807v2 (CC BY 4.0)}
\begin{document}
\maketitle

\noindent\emph{Source and licence.} A new bound for the Kalton--Roberts constant. Version arXiv:2606.06807v2 (CC BY 4.0), distributed by arXiv under the Creative Commons Attribution 4.0 International licence, http://creativecommons.org/licenses/by/4.0/, as recorded in the arXiv licence metadata for this exact version and shown on the abstract page https://arxiv.org/abs/2606.06807v2. Full licence notice: \texttt{licenses/arxiv-2606.06807-CC-BY-4.0.txt}.\par\smallskip

\noindent\textit{Editorial scope.} Counted unit: the article's lemma lem:extremal-families (source0.tex lines 238--246). The article's complete original proof of it is delivered with it (source0.tex lines 248--278, one \texttt{proof} environment of 31 lines). No mathematics is written, repaired or re-derived: the statement and the proof are the article's own lines, in the article's own notation, and no retained line is substituted or re-flowed.  No further source-local context is needed: every cross-reference of every retained line resolves inside the two delivered spans.  Nothing was omitted from any retained span, so the package carries no omission record.  No retained line contains a citation, so the wrapper carries no bibliography and no bibliography entry.  No retained line contains a figure environment, an image-inclusion command or drawing code, so no image is delivered and the package declares no asset dependency.  Editorial formatting only, in addition: an AMS \texttt{amsart} preamble that restates the article's own declarations and macros used by the retained lines, namely the environments \texttt{lemma} on separate counters, together with the standard packages those lines need.  The retained environments print in this document as Lemma 1 in order of appearance, whereas the article prints the same items with the numbering of its own full document.  The complete native source of the article as distributed in the arXiv e-print for this exact version is bundled unchanged as \texttt{sources/202213/source0.tex}.
\begin{lemma}\label{lem:extremal-families}
There are rational weighted families \(\cA\) and \(\cB\) such that
\begin{enumerate}[label=\textup{(\roman*)}]
\item every point has frequency \(q\) in \(\cA\), and the average of
      \(M-f(A)\) over \(\cA\) is at most \(q(1-u)\);
\item every point has frequency \(p\) in \(\cB\), and the average of
      \(M+f(B)\) over \(\cB\) is at most \(p(1+u)\).
\end{enumerate}
\end{lemma}

\begin{proof}
Write \(\lambda=\lambda^+-\lambda^-\), and for \(i\in U\) put
\[
 m_i^+=\sum_{\substack{A\ni i}}\lambda_A^+,
 \qquad
 m_i^-=\sum_{\substack{A\ni i}}\lambda_A^-.
\]
The zero-marginal equations say that \(m_i^+=m_i^-\).

Form \(\cA\) from the positive active sets with their \(\lambda^+\)-weights
and the complements of the negative active sets with their
\(\lambda^-\)-weights.  Its total weight is one, and the frequency of \(i\)
is
\[
        m_i^++(q-m_i^-)=q.
\]
Positive active sets have deficit zero.  If \(A_-\) is negative active, then
\(A_-\sqcup A_-^c=U\), and \(1\)-additivity gives
\[
        M-f(A_-^c)\leqslant1-u.
\]
Averaging gives (i).  Define \(\cB\) using the negative active sets with
their \(\lambda^-\)-weights and the complements of the positive active sets
with their \(\lambda^+\)-weights.  The same calculation gives point
frequency \(p\).
The inequality
\[
        M+f(A_+^c)\leqslant1+u
\]
for positive active \(A_+\) gives (ii).
\end{proof}

\end{document}
